\section{The starting point changes the sign}
\label{sec:starting}

The two benchmarks differ in a way that is usually treated as a nuisance to be
disclosed. TabReD supplies nothing and the task is solved from scratch;
MLAgentBench supplies a working training script and asks for improvement. We
treat this as a factor rather than a caveat, because the value of prescriptive
structure depends on it.

The win--loss counts split cleanly along that line. In the greenfield
condition---the eight TabReD tasks, where nothing is supplied beyond the task
statement---the branch carrying the prescription wins all eight, at three
attempts per cell on the common adapter (Table~\ref{tab:grid}, column
\emph{both}). In the improve-a-baseline condition---the MLAgentBench tasks,
which each ship a working \texttt{train.py}---the branch without it wins five of
the eight we were able to run three times, and three of those five are total:
the prescribed branch submits nothing at all in three attempts
(Table~\ref{tab:ktool}). Both counts come from
the same system, the same backbone and the same pair of instruction files; only
the starting point differs. The reading is the one the mechanism predicts: where
there is no default to override, any prescription that points the model at a
reasonable tool improves on what it would otherwise construct, and where the
supplied script already contains the decision the prescription was introduced to
make, overriding it discards work that was correct.

The counts are symmetric but the magnitudes are not, and that asymmetry is the
more informative half. Improving a baseline, the prescription can take a task
from a usable score to no submission at all, three attempts of three
(Table~\ref{tab:ktool}). On greenfield nothing of that kind happens anywhere:
every cell returns a metric on every task, and the normalised effect of the
layer runs from $+0.004$ on \texttt{sberbank-housing} to $+0.357$ on
\texttt{ecom-offers} with a median of $0.06$. Six of the eight exceed $0.020$,
the largest run-to-run range anywhere in that arm; the two that do not, $+0.004$
and $+0.005$, fall on tasks whose attempts reproduce bit-identically, so their
apparent spread of zero measures library determinism rather than agent stability
(Section~\ref{sec:design}) and we rest nothing on them. Even taken at face
value the effect is modest on most tasks, for the reason
Section~\ref{sec:spontaneous} gives: the system reaches for the prescribed
family unaided. Greenfield does not mean structure pays richly; it means
structure stops costing.

The two halves of this comparison are not a crossed design, and the aggregate
win--loss count across them should not be read as a single number. The
benchmarks differ in their tasks, metrics and time budgets as well as in their
starting point, so what we observe is one factor across two protocols. Closing
that confound needs either a greenfield variant of the MLAgentBench tasks or an
improve-a-baseline variant of the TabReD ones, built so that only the starting
point moves; we did not build one, and state the weaker claim accordingly. What
the evidence does support without the crossing is the asymmetry itself, which is
visible inside each protocol separately: prescription never destroys a task where
nothing was supplied, and destroys three of eight where a working script was.
